% =====================================================================
% Generative AI - Assignment 2 - Springer LNCS Report
% Muhammad Ahmed Mufti (22I-1088)
% =====================================================================
\documentclass[runningheads]{llncs}

\usepackage[T1]{fontenc}
\usepackage[utf8]{inputenc}
\usepackage{graphicx}
\usepackage{amsmath,amssymb}
\usepackage{booktabs}
\usepackage{hyperref}
\usepackage{xcolor}
\usepackage{caption}
\usepackage{float}
\usepackage{tabularx}
\usepackage{array}

\hypersetup{colorlinks=true, linkcolor=black, citecolor=black, urlcolor=blue}
\graphicspath{{images/}}

\begin{document}

% =====================================================================
\title{Generative AI, Assignment~2:\\
CycleGAN Image Translation, English--Urdu Machine Translation with
Transformers, and Vision Transformer vs.\ CNN on CIFAR-10}
\titlerunning{GenAI Assignment 2}

\author{Muhammad Ahmed Mufti\inst{1}}
\authorrunning{M.\ A.\ Mufti}

\institute{Department of Computer Science,\\
National University of Computer and Emerging Sciences,\\
Islamabad, Pakistan\\
\email{i221088@nu.edu.pk}}

\maketitle

% =====================================================================
\begin{abstract}
We implement and evaluate three foundational deep-learning architectures
for generative and discriminative tasks. First, a CycleGAN is trained
end-to-end on the Person Face Sketches dataset for unpaired
image-to-image translation between photographs and pencil sketches, and
is deployed behind a Flask web interface that supports file upload,
webcam capture and automatic input-type detection. Second, a Transformer
encoder--decoder is implemented from scratch for English-to-Urdu neural
machine translation, trained with Byte-Pair Encoding subword
tokenisation, Noam learning-rate scheduling, label smoothing and EDA
data augmentation, and compared against fine-tuned mBART-50. Third, a
ResNet-style CNN, a from-scratch Vision Transformer and a pretrained
ViT-tiny are compared on CIFAR-10 across accuracy, F1-score, confusion
matrices, training time and parameter count. All training runs use
per-epoch checkpointing on a persistent Modal volume and are resumable
after interruption. We report quantitative results, qualitative
examples, and discuss the architectural trade-offs empirically.

\keywords{CycleGAN \and Transformer \and Machine Translation \and
Vision Transformer \and CNN \and CIFAR-10 \and Transfer Learning \and
Data Augmentation.}
\end{abstract}

% =====================================================================
\section{Introduction}

This report documents three interconnected generative-AI experiments.
Task~1 trains an unpaired image-to-image translation network
(CycleGAN~\cite{cyclegan}) between real face photographs and pencil
sketches on the Person Face Sketches
dataset~\cite{personfacesketches}, and provides a Flask-based user
interface for live inference. Task~2 develops an English-to-Urdu
machine translation system using the Transformer encoder--decoder of
Vaswani~et~al.~\cite{vaswani} implemented from scratch in PyTorch, and
compares it against both zero-shot and fine-tuned mBART-50
baselines~\cite{mbart}. Task~3 contrasts a custom CNN baseline against
a Vision Transformer~\cite{vit} trained from scratch and a pretrained
ViT-tiny fine-tuned on CIFAR-10~\cite{cifar10}. The three tasks share a
common infrastructure: every training loop saves a per-epoch checkpoint
to a persistent Modal Volume and supports \texttt{--resume}, so that
sessions can be interrupted and resumed without loss of progress.

The remainder of the paper is organised as follows. Sections~2--4
describe each task in turn, with subsections covering method,
implementation, training procedure, quantitative results and
qualitative analysis. Section~5 aggregates findings and discusses
cross-task observations; Section~6 concludes.

% =====================================================================
\section{Task 1: CycleGAN for Face $\leftrightarrow$ Sketch Translation}

\subsection{Dataset and Pre-processing}

We use the Person Face Sketches Kaggle dataset~\cite{personfacesketches},
which contains aligned photograph/sketch pairs. We use the provided
\texttt{train/} split (ignoring \texttt{val/}, \texttt{test/} at
training time) with the photographs forming domain~$A$ and the pencil
sketches domain~$B$. Although the dataset is nominally paired, CycleGAN
does not require correspondence, so we sample from each domain
\emph{independently} at training time. This mirrors the original
paper's unpaired setting and guards the model from learning any
spurious pair-level shortcut.

Each image is resized so that its short side is $286$ pixels, randomly
cropped to $256 \times 256$, horizontally flipped with probability 0.5,
and normalised to $[-1, 1]$ per channel. At inference and for the UI
sample grid the test-time pipeline is a single $256 \times 256$
resize and normalise, without augmentation.

\subsection{Model Architecture}

Two generators $G_{AB}: A \to B$ and $G_{BA}: B \to A$ and two
discriminators $D_A, D_B$ are trained adversarially. Each generator is
the ResNet-based architecture from~\cite{cyclegan}:

\begin{itemize}
\item \textbf{Initial block:} $\mathrm{ReflectionPad}(3)$ +
$\mathrm{Conv}(3\!\to\!64, 7\times 7)$ + InstanceNorm + ReLU.
\item \textbf{Two strided downsampling blocks:} $\mathrm{Conv}(64\!\to\!128)$
and $\mathrm{Conv}(128\!\to\!256)$, each with InstanceNorm + ReLU and
stride 2.
\item \textbf{Nine residual blocks} at 256 channels with reflection
padding and two $3\times 3$ convolutions each.
\item \textbf{Two nearest-neighbour upsampling blocks} back to 128 and
64 channels.
\item \textbf{Final block:} $\mathrm{ReflectionPad}(3)$ +
$\mathrm{Conv}(64\!\to\!3, 7\times 7)$ + $\tanh$.
\end{itemize}

The discriminators follow the $70\times 70$ PatchGAN design: four
strided $4\times 4$ convolutions with LeakyReLU(0.2), producing a
$30\times 30$ logit map that scores each $70\times 70$ input patch as
real or fake. All weights are initialised Normal$(0, 0.02)$, matching
the paper.

\subsection{Losses}

CycleGAN combines four loss terms. Using LSGAN~\cite{lsgan} rather than
vanilla BCE for the adversarial objective:
\begin{align}
\mathcal{L}_{\mathrm{GAN}}(G, D, X, Y) &=
\mathbb{E}_{y\sim Y}\!\left[(D(y)-1)^2\right]
+ \mathbb{E}_{x\sim X}\!\left[D(G(x))^2\right], \\
\mathcal{L}_{\mathrm{cyc}}(G_{AB}, G_{BA}) &=
\mathbb{E}_a\!\left[\lVert G_{BA}(G_{AB}(a)) - a \rVert_1\right]
+ \mathbb{E}_b\!\left[\lVert G_{AB}(G_{BA}(b)) - b \rVert_1\right], \\
\mathcal{L}_{\mathrm{id}}(G_{AB}, G_{BA}) &=
\mathbb{E}_a\!\left[\lVert G_{BA}(a) - a \rVert_1\right]
+ \mathbb{E}_b\!\left[\lVert G_{AB}(b) - b \rVert_1\right].
\end{align}
The total generator objective is
\begin{equation}
\mathcal{L}_G = \mathcal{L}_{\mathrm{GAN}} +
\lambda_{\mathrm{cyc}}\,\mathcal{L}_{\mathrm{cyc}} +
\lambda_{\mathrm{id}}\,\mathcal{L}_{\mathrm{id}},
\end{equation}
with $\lambda_{\mathrm{cyc}}=10$ and $\lambda_{\mathrm{id}}=5$. We also
maintain a 50-image replay buffer for each discriminator update, as
prescribed in~\cite{cyclegan}, to stabilise training by exposing $D$
to a history of previous fake samples rather than only the current
batch.

\subsection{Training procedure}

Training was run in two phases: an initial 5~epoch run at $128\times
128$ with batch size~8 on a Colab T4, followed by 30~epochs at the
paper-faithful $256\times 256$ batch~8 on a Modal serverless L4 GPU.
Optimiser: Adam with $\beta_1=0.5$, $\beta_2=0.999$,
$\eta=2\!\times\!10^{-4}$. A piecewise-linear learning-rate schedule
keeps $\eta$ constant for the first half of training and then linearly
decays it to zero over the remaining half. After every epoch the
script dumps a full checkpoint (generator, discriminator and optimiser
states) to a persistent volume and writes a generator-only file
suitable for inference. A sample grid showing four real photos, their
generated sketches, four real sketches and their generated faces is
saved as PNG each epoch for visual monitoring.

\paragraph{Checkpoint selection.} Empirically we observed that the two
translation directions converge at different rates. With batch size 8
(larger than the paper's batch 1), the $G_{AB}$ sketch-generation head
continues to sharpen its output through all 30 epochs, while the
$G_{BA}$ face-reconstruction head reaches its visual-quality peak
around epoch 5 and subsequently deteriorates --- the generator trades
photographic fidelity for artistic texture as cycle and adversarial
losses push it away from the real-photo manifold. We therefore select
the early ($128\times 128$, 5-epoch) checkpoint for the final
deployed generators, applying the standard early-stopping heuristic
based on qualitative visual inspection of the per-epoch sample grids.
This observation is itself a useful finding: visually inspecting the
sample grid at every epoch is preferable to blindly training for the
paper's full 100--200 epochs on a dataset of this size.

\paragraph{Stability techniques.} We use InstanceNorm instead of
BatchNorm throughout because the normalisation is performed per-sample,
which avoids mixing style statistics across the batch (a common source
of artefacts in GAN-based style transfer). Reflection padding around
each convolution prevents the visible edge halos that zero padding
produces. The replay buffer softens discriminator updates and the
piecewise-linear LR decay prevents the generators from diverging late
in training. Finally, because the original 256~px batch-1 training
from the paper would require 10+~hours on a T4, we cautiously scale the
batch size to 8 while keeping everything else identical; this matches
our memory budget on L4 without harming output quality.

\paragraph{Checkpointing and resume.} The training script writes
\texttt{weights/latest.pt}, \texttt{weights/epoch\_NNN.pt} and
\texttt{weights/generators.pt} after every epoch. Passing
\texttt{--resume} on the next launch loads \texttt{latest.pt} and
continues from the saved epoch, optimiser state included. The Modal
Volume that backs these files persists across container restarts, so
even if the training job is interrupted (e.g.\ platform timeout,
network disconnect) training resumes without loss.

\subsection{Web Interface for Real-Time Inference}

A Flask web application (\texttt{flask\_app.py}) wraps the trained
generators behind a browser UI (Fig.~\ref{fig:q1_ui_face},
Fig.~\ref{fig:q1_ui_sketch}) with three input modes:

\begin{itemize}
\item \textbf{File upload} --- any JPEG/PNG, arbitrary resolution.
\item \textbf{Webcam capture} --- using
\texttt{navigator.mediaDevices.getUserMedia} in the front-end, a
dataURL of the captured frame is POSTed to the \texttt{/translate}
endpoint.
\item \textbf{Manual override} --- user may force face-to-sketch or
sketch-to-face mode to bypass detection.
\end{itemize}

An automatic input-type detector classifies the image as sketch or
photograph by measuring per-pixel deviation from the grayscale average:
if $\mathrm{mean}_{i}\lVert x_i - \bar{x}_i \rVert_1 < 12$ (on a 0--255
scale), the image is treated as a sketch and routed through $G_{BA}$;
otherwise it is routed through $G_{AB}$. The UI reports both the
chosen mode and the raw detector value for transparency, and exposes a
\texttt{/status} endpoint to confirm the model has loaded.

\subsection{Results}

Figure~\ref{fig:q1_samples} shows a late-epoch sample grid: row~1 real
photographs, row~2 their generated sketches, row~3 real sketches, row~4
their generated faces. The generated sketches are clean pencil-like
line drawings that preserve identity and pose; the generated faces
reconstruct recognisable human features even from sparse inputs,
although fine details such as lip colour and eye highlights are still
blurry at 30~epochs. Adversarial, cycle and identity losses all decreased
monotonically through training. Figures~\ref{fig:q1_ui_face} and
\ref{fig:q1_ui_sketch} show the Flask UI rendering translations of
novel uploaded images.

\begin{figure}[H]
\centering
\includegraphics[width=0.9\textwidth]{Q1_CycleGAN/03_sample_epoch_final.png}
\caption{CycleGAN sample grid from the selected checkpoint. Row 1:
real photographs (domain~$A$). Row~2: generated sketches $G_{AB}(A)$.
Row~3: real sketches (domain~$B$). Row~4: generated photographs
$G_{BA}(B)$.}
\label{fig:q1_samples}
\end{figure}

\begin{figure}[H]
\centering
\includegraphics[width=0.85\textwidth]{Q1_CycleGAN/04_ui_face_to_sketch.png}
\caption{Flask web interface: uploaded face photographs (left) are
automatically routed through $G_{AB}$ and rendered as pencil sketches
(right).}
\label{fig:q1_ui_face}
\end{figure}

\begin{figure}[H]
\centering
\includegraphics[width=0.85\textwidth]{Q1_CycleGAN/05_ui_sketch_to_face.png}
\caption{Flask web interface: uploaded sketches (left) are routed
through $G_{BA}$ and rendered as reconstructed faces (right).}
\label{fig:q1_ui_sketch}
\end{figure}

\subsection{Final Evaluation and Discussion}

The selected checkpoint achieves two paper-faithful design hallmarks:
(a) outputs are semantically aligned with inputs (pose, hairstyle,
facial structure preserved) and (b) cycle consistency prevents mode
collapse. The generated sketches are clean pencil-style line drawings
that retain identity; the generated faces have recognisable features
with natural colour and lighting. A side observation from training was
that the two directions optimise at different rates: longer training
helps sketch quality but degrades face reconstruction. This motivates
either bidirectional early stopping or asymmetric training objectives
in future work. The Flask UI handles both directions of the
translation without code duplication and runs at interactive speed on
CPU ($\approx$2~s per 256~px image), making it suitable for live demos.

% =====================================================================
\section{Task 2: English-to-Urdu Transformer Machine Translation}

\subsection{Dataset and Pre-processing}

We use the \texttt{parallel-corpus-for-english-urdu-language} Kaggle
dataset~\cite{zainparallel}, which contains 24\,524 aligned
English-Urdu sentence pairs as two line-aligned text files.
Figure~\ref{fig:q2_align} demonstrates the alignment at the first three
lines. We split the corpus 93/5/2 into train / validation / test
(22\,808 / 1\,226 / 490 pairs) using a fixed seed for reproducibility.
All sentences are normalised by trimming whitespace and collapsing
runs of internal whitespace.

\begin{figure}[H]
\centering
\includegraphics[width=0.9\textwidth]{Q2_Translation/02_corpus_alignment.png}
\caption{Line-level alignment check: the first three lines of the
English and Urdu corpora correspond semantically.}
\label{fig:q2_align}
\end{figure}

\subsection{Subword Tokenisation}

We train two separate Byte-Pair Encoding (BPE) tokenisers, one per
language, using the HuggingFace \texttt{tokenizers} library.
Vocabulary size is 8\,000 per side; the special tokens \texttt{<pad>},
\texttt{<bos>}, \texttt{<eos>}, \texttt{<unk>} occupy the first four
IDs. BPE handles Urdu's rich morphology and produces a compact
vocabulary: on the training split we obtain 7\,957 effective English
sub-tokens and 7\,496 Urdu sub-tokens. At encoding time we prepend
\texttt{<bos>} and append \texttt{<eos>}; decoding strips these
markers.

\subsection{Data Augmentation (EDA)}

Following Wei and Zou~\cite{eda}, we implement three source-side
augmentation operators (augmentation is applied only to the English
input; the Urdu target remains faithful to its label):

\begin{itemize}
\item \textbf{Random deletion} --- each whitespace-separated token is
dropped independently with probability~0.1.
\item \textbf{Random swap} --- swap a randomly chosen pair of adjacent
tokens.
\item \textbf{Synonym replacement} --- replace one token with a
WordNet synonym (NLTK).
\end{itemize}

At training time each sample is augmented with probability
$p_{\mathrm{aug}} = 0.4$ by selecting one operator uniformly at random.
Validation and test inputs are never augmented.
Figure~\ref{fig:q2_aug} shows representative augmented outputs.

\begin{figure}[H]
\centering
\includegraphics[width=0.9\textwidth]{Q2_Translation/07_data_augmentation_samples.png}
\caption{Source-side EDA augmentation: for each English sentence the
first line is the original, the following three are random augmented
variants.}
\label{fig:q2_aug}
\end{figure}

\subsection{Model Architecture}

The Transformer encoder--decoder of~\cite{vaswani} is implemented from
scratch in PyTorch (\texttt{transformer\_model.py}); we \emph{do not}
use \texttt{nn.Transformer}. Every sub-module is explicit:

\begin{itemize}
\item \textbf{Embeddings.} Separate source- and target-side
\texttt{nn.Embedding} layers of width $d_{\mathrm{model}}{=}256$,
scaled by $\sqrt{d_{\mathrm{model}}}$ before addition to the positional
encoding.
\item \textbf{Sinusoidal positional encoding.} $PE_{(pos, 2i)} =
\sin(pos/10000^{2i/d_{\mathrm{model}}})$,
$PE_{(pos, 2i+1)} = \cos(pos/10000^{2i/d_{\mathrm{model}}})$.
\item \textbf{Multi-head scaled dot-product attention.} Eight heads of
dimension 32 each, with independent Q/K/V linear projections and a
final output projection. A broadcastable mask nullifies attention to
padding positions in the encoder and to future positions in the
decoder self-attention.
\item \textbf{Position-wise feed-forward} sub-layer of width $d_{\mathrm{ff}} = 1024$
with ReLU activation and dropout 0.1.
\item \textbf{Residual connections} around each sub-layer with
post-norm LayerNorm.
\item \textbf{Final projection} to target vocabulary with no weight
tying.
\end{itemize}

Both encoder and decoder have 4 stacked layers, yielding 13.25M
trainable parameters in total. Figure~\ref{fig:q2_header} shows the
training-time model summary printed by the script.

\begin{figure}[H]
\centering
\includegraphics[width=0.95\textwidth]{Q2_Translation/01_training_header.png}
\caption{Dataset and model configuration summary at the start of
training.}
\label{fig:q2_header}
\end{figure}

\subsection{Training Process}

We use Adam with $\beta_1=0.9,\beta_2=0.98,\epsilon=10^{-9}$ as in the
original paper. The Noam learning-rate schedule is
\begin{equation}
\eta(s) = d_{\mathrm{model}}^{-0.5}\cdot
\min\!\left(s^{-0.5},\;s\cdot w^{-1.5}\right),
\end{equation}
with warm-up steps $w=2000$. Loss is label-smoothed cross-entropy
\begin{equation}
\mathcal{L}_{\mathrm{ls}} = (1-\varepsilon)\,\mathrm{NLL}(p, y) +
\varepsilon\,\mathbb{E}_{v}\!\left[-\log p(v)\right],
\quad \varepsilon=0.1,
\end{equation}
which discourages overconfidence and empirically improves BLEU.
Gradient clipping ($\lVert g\rVert_2 \le 1$) is applied at every step
to stabilise early training. Batch size is 64 sentences, maximum
sequence length 128 sub-tokens. Training ran on a Modal L4 GPU with
per-epoch checkpointing and \texttt{--resume} support. We trained for
20 initial epochs without augmentation, then 20 additional epochs with
$p_{\mathrm{aug}} = 0.4$ augmentation enabled.

\subsection{Evaluation}

\paragraph{BLEU.} After each epoch the model is evaluated on a
200-sample subset of the validation set with greedy decoding, using
\texttt{sacrebleu}'s corpus-level BLEU implementation. Pre-augmentation
final BLEU was 0.53; post-augmentation final BLEU was 0.43 (40~epochs
total). Qualitative inference (Fig.~\ref{fig:q2_collapse}) confirms
what the low BLEU suggests: the greedy decoder converges to a single
high-frequency pattern, producing identical Urdu output regardless of
input --- the classical ``mode collapse'' failure mode observed in
vanilla Transformer MT trained on small corpora.

\begin{figure}[H]
\centering
\includegraphics[width=0.9\textwidth]{Q2_Translation/04_from_scratch_mode_collapse.png}
\caption{Inference from the from-scratch Transformer on five test
sentences: the output is identical for every input, although the
target references vary substantially.}
\label{fig:q2_collapse}
\end{figure}

\paragraph{Analysis.} The result is consistent with the literature:
Vaswani et al.~\cite{vaswani} trained on $\approx$4.5M WMT14 EN-DE
pairs---roughly $200\times$ our corpus size. Our 22\,808 pairs are
below the threshold at which a vanilla Transformer with $\sim$13M
parameters can avoid the collapse-to-prior failure mode. Data
augmentation introduced lexical variation in the source (verified by
manual inspection) but did not fundamentally bridge the
information-theoretic gap on the target side, so the collapse persists.

\subsection{Pre-trained Baselines: Zero-Shot and Fine-Tuned mBART-50}

We compare the from-scratch model against the publicly available
\texttt{facebook/mbart-large-50-many-to-many-mmt}~\cite{mbart}, a
610M-parameter multilingual model pretrained on 50 languages including
Urdu. The zero-shot baseline produces fluent Urdu translations for
four of our five test sentences without any in-domain fine-tuning
(Fig.~\ref{fig:q2_mbart}); the fifth (``Education is the key to
success.'') is near-correct but transliterates the English word
``education''. We then fine-tune mBART-50 for three epochs on our
22\,808 pairs (batch 4, AdamW, $\eta=3\!\times\!10^{-5}$, 10\%
warm-up, label smoothing inherited from mBART's default) on a Modal
A100-40GB; final translations in Fig.~\ref{fig:q2_mbart_ft} are
noticeably more idiomatic (for example the second sentence switches
from ``main'' to the more natural ``mujhe'') and the ``edukation''
transliteration is replaced by the correct Urdu word \emph{ta`leem}.

\begin{figure}[H]
\centering
\includegraphics[width=0.9\textwidth]{Q2_Translation/05_mbart_zeroshot_translations.png}
\caption{Zero-shot mBART-50 translations (no fine-tuning).}
\label{fig:q2_mbart}
\end{figure}

\begin{figure}[H]
\centering
\includegraphics[width=0.9\textwidth]{Q2_Translation/08_mbart_finetuned_translations.png}
\caption{mBART-50 after three epochs of fine-tuning on our corpus.}
\label{fig:q2_mbart_ft}
\end{figure}

\begin{figure}[H]
\centering
\includegraphics[width=\textwidth]{Q2_Translation/06_comparison_card.png}
\caption{Side-by-side qualitative comparison of the from-scratch
Transformer (collapsed) against zero-shot mBART-50 on the same five
evaluation sentences.}
\label{fig:q2_compare}
\end{figure}

\subsection{Final Performance Summary}

Table~\ref{tab:q2_summary} summarises Q2. The clear lesson is that at
the 22\,808-pair scale, from-scratch Transformer MT is insufficient,
while large-scale multilingual pretraining (mBART) transfers
essentially for free. EDA augmentation on this dataset did not
compensate for the scale deficit, although it is still a valuable
technique at larger corpus sizes.

\begin{table}[H]
\centering
\caption{Q2 model summary.}
\label{tab:q2_summary}
\renewcommand{\arraystretch}{1.2}
\begin{tabular}{lrcl}
\toprule
Model & Parameters & BLEU (val, 200) & Qualitative \\
\midrule
Transformer from scratch, 20~ep                        & 13.3 M  & 0.53 & collapsed \\
Transformer from scratch + EDA augmentation, 40~ep     & 13.3 M  & 0.43 & collapsed (different phrase) \\
mBART-50 zero-shot                                     & 610 M   & n/a  & 4/5 correct \\
mBART-50 fine-tuned, 3~ep                              & 610 M   & n/a  & 4/5 correct, more idiomatic \\
\bottomrule
\end{tabular}
\end{table}

% =====================================================================
\section{Task 3: Vision Transformer vs.\ CNN on CIFAR-10}

\subsection{Dataset and Pre-processing}

The CIFAR-10 dataset~\cite{cifar10} contains 60\,000 $32\times 32$
colour images across ten balanced classes. We use the canonical 50\,000
/ 10\,000 train/test split provided by \texttt{torchvision}. Training
augmentation: random crop with 4-pixel reflection padding, horizontal
flip. Normalisation uses the CIFAR-10 channel means
$(0.491, 0.482, 0.447)$ and standard deviations
$(0.247, 0.244, 0.262)$. Evaluation uses only normalisation.

\subsection{Model Architectures}

\paragraph{CNN baseline.} A ResNet-20-style convolutional network with
a 64-channel stem followed by three stages of two \textsc{BasicBlock}
residual units each (64$\to$128$\to$256 channels, strided downsampling
between stages), global average pooling, dropout 0.3 and a linear
classification head. Each residual block contains two $3\times 3$
convolutions with batch-normalisation and ReLU. Total: 2\,777\,674
trainable parameters.

\paragraph{Vision Transformer from scratch.} Following Dosovitskiy
et~al.~\cite{vit}: $4\times 4$ patches on a $32\times 32$ input yield
$8\times 8 = 64$ tokens. The token embedding is produced by a
$\mathrm{Conv2d}(3, d_{\mathrm{model}}, \mathrm{kernel}=4, \mathrm{stride}=4)$
which is equivalent to flatten-then-linear. A learnable \texttt{[CLS]}
token is prepended and a learnable 65-entry positional embedding is
added. Six pre-norm transformer blocks follow, each with multi-head
self-attention ($h=6$ heads of dimension 32) and an MLP of ratio 4
with GELU activation. Final LayerNorm on the \texttt{[CLS]} token feeds
a linear classification head. Total: 2\,693\,578 trainable parameters.
Parameter budget is intentionally matched to the CNN for a fair
comparison.

\paragraph{Pretrained ViT-tiny.} We fine-tune
\texttt{vit\_tiny\_patch16\_224} from the \texttt{timm}
library~\cite{timm}, pretrained on ImageNet-21k. Inputs are upsampled
to $224\times 224$ to match the pretraining resolution. The model has
5\,526\,346 parameters.

\subsection{Training Procedure}

All three models share the same training loop (\texttt{train.py}):
AdamW with weight decay $5\!\times\!10^{-4}$, cosine learning-rate
schedule over the full epoch budget, label smoothing 0.1 on the
cross-entropy loss, gradient clipping $\lVert g\rVert_2 \le 1$, and
early stopping on validation accuracy with patience 10. Per-epoch
checkpointing writes \texttt{latest.pt} and \texttt{best.pt} (the
model with highest val accuracy so far); \texttt{--resume} loads
\texttt{latest.pt}.

Hyperparameter choices are summarised in Table~\ref{tab:q3_hp}. Batch
sizes were chosen to maximise GPU utilisation at the given image
resolution. Learning rates were selected by a small preliminary sweep
($\{10^{-4}, 3\!\times\!10^{-4}, 5\!\times\!10^{-4}, 10^{-3}\}$) on a
10\% subset.

\begin{table}[H]
\centering
\caption{Q3 hyperparameters per model.}
\label{tab:q3_hp}
\renewcommand{\arraystretch}{1.2}
\begin{tabular}{lccccc}
\toprule
Model & Image & Batch & LR & Epochs & Augmentation \\
\midrule
CNN                  & $32\!\times\!32$  & 128 & $3\!\times\!10^{-4}$ & 50 & crop + flip \\
ViT from scratch     & $32\!\times\!32$  & 128 & $5\!\times\!10^{-4}$ & 80 (ES at 55) & crop + flip \\
Pretrained ViT-tiny  & $224\!\times\!224$ & 64 & $1\!\times\!10^{-4}$ & 3  & crop + flip \\
\bottomrule
\end{tabular}
\end{table}

\subsection{Evaluation Metrics}

For each model we record per-epoch training and validation loss +
accuracy. On the best checkpoint we additionally compute macro-averaged
precision, recall and F1 using \texttt{scikit-learn}, a $10\times 10$
confusion matrix (normalised), a classification report, and a 16-image
example-prediction grid. All metrics are dumped to
\texttt{results.json} alongside the plots.

\subsection{Quantitative Results}

Table~\ref{tab:q3} summarises the final metrics, and
Fig.~\ref{fig:q3_compare} provides a visual comparison.
Fig.~\ref{fig:q3_overlay} shows the three validation-accuracy curves on
a shared axis; Fig.~\ref{fig:q3_curves} the CNN's and ViT's individual
loss and accuracy trajectories; Fig.~\ref{fig:q3_pretrained} the
pretrained ViT convergence.

\begin{table}[H]
\centering
\caption{CIFAR-10 validation metrics after full training.}
\label{tab:q3}
\renewcommand{\arraystretch}{1.2}
\begin{tabular}{lrrrrr}
\toprule
Model & Params & Val.\ acc & F1 (macro) & Time & Epochs \\
\midrule
CNN (ResNet-20-like)  & 2\,777\,674 & 0.9278 & 0.928 & 31 min & 50 \\
ViT (from scratch)    & 2\,693\,578 & 0.8219 & 0.821 & 36 min & 55 (ES) \\
Pretrained ViT-tiny   & 5\,526\,346 & \textbf{0.9715} & \textbf{0.972} & 10 min & 3 \\
\bottomrule
\end{tabular}
\end{table}

\begin{figure}[H]
\centering
\includegraphics[width=0.48\textwidth]{Q3_ViT_CNN/01_cnn_curves.png}
\hfill
\includegraphics[width=0.48\textwidth]{Q3_ViT_CNN/02_vit_curves.png}
\caption{Loss and accuracy curves: CNN (left), ViT-from-scratch
(right). The ViT converges more slowly and plateaus at lower
validation accuracy.}
\label{fig:q3_curves}
\end{figure}

\begin{figure}[H]
\centering
\includegraphics[width=0.6\textwidth]{Q3_ViT_CNN/03_pretrained_curves.png}
\caption{Pretrained ViT-tiny: val accuracy exceeds 93\% after the first
epoch and stabilises at $\sim$97\%.}
\label{fig:q3_pretrained}
\end{figure}

\begin{figure}[H]
\centering
\includegraphics[width=0.95\textwidth]{Q3_ViT_CNN/08_comparison_chart.png}
\caption{Accuracy and parameter count across the three models.}
\label{fig:q3_compare}
\end{figure}

\begin{figure}[H]
\centering
\includegraphics[width=0.9\textwidth]{Q3_ViT_CNN/09_val_acc_overlay.png}
\caption{Validation accuracy per epoch overlaid for all three models.
The pretrained ViT immediately dominates; the from-scratch ViT takes
many epochs to begin to close the gap with the CNN.}
\label{fig:q3_overlay}
\end{figure}

\subsection{Per-Class Analysis}

The confusion matrices in Fig.~\ref{fig:q3_cm} reveal consistent
patterns: the CNN's error mass concentrates on \textit{cat}/\textit{dog},
\textit{deer}/\textit{horse} and \textit{bird}/\textit{airplane}
confusions. The from-scratch ViT shows substantially more diffuse
error, with higher confusion among all living-creature classes and a
notable weakness on \textit{bird} and \textit{cat}. The pretrained ViT
has a much tighter diagonal, with the only visible residual confusion
between the visually-similar \textit{cat}/\textit{dog} pair.

\begin{figure}[H]
\centering
\includegraphics[width=0.32\textwidth]{Q3_ViT_CNN/04_cnn_confusion.png}
\hfill
\includegraphics[width=0.32\textwidth]{Q3_ViT_CNN/05_vit_confusion.png}
\hfill
\includegraphics[width=0.32\textwidth]{Q3_ViT_CNN/06_pretrained_confusion.png}
\caption{Normalised confusion matrices: CNN (left),
ViT-from-scratch (centre), pretrained ViT-tiny (right).}
\label{fig:q3_cm}
\end{figure}

\begin{figure}[H]
\centering
\includegraphics[width=0.6\textwidth]{Q3_ViT_CNN/07_example_preds.png}
\caption{Example predictions: green titles indicate correct
classification, red indicate errors.}
\label{fig:q3_examples}
\end{figure}

\subsection{Attention vs.\ Spatial Modelling}

The experiment isolates the practical consequences of two very
different inductive biases:

\begin{itemize}
\item \textbf{CNNs} enforce local receptive fields and translation
equivariance by construction. At CIFAR's $32\times 32$ resolution
these priors essentially encode the correct shape of the function that
maps patches to object classes.
\item \textbf{ViTs} trade those priors for content-dependent global
connectivity through self-attention. Every token can attend to every
other token from the first layer, which is extremely expressive but
informationally underconstrained when data is scarce.
\end{itemize}

Empirically the consequences are clear: at matched $\approx$2.7M
parameters, the CNN outperforms the from-scratch ViT by 10.6 absolute
accuracy points on CIFAR-10. This is not a failure of the ViT
architecture --- it is a direct confirmation of the original paper's
result that ViT beats CNNs only at ImageNet scale, because only there
is the inductive prior of locality/equivariance learnable from data.
The pretrained ViT, which was exposed to ImageNet-21k during its
supervised pretraining, has effectively already learnt these priors;
fine-tuning for three epochs then matches them to CIFAR's statistics
and we observe a +4.4 point advantage over the best CNN result with
one third of the training time. In other words, attention is not
\emph{worse} than convolution; it just requires either much more data
or transfer from a related source to outperform it.

\subsection{Final Report and Interpretation}

The three-way comparison provides a compact narrative of why modern
computer vision converged on transfer learning: (i) under a fixed
compute and data budget a well-designed CNN remains the strong
baseline; (ii) a matched-capacity transformer underperforms this
baseline; and (iii) pretraining + lightweight fine-tuning flips the
ordering and reaches the highest absolute accuracy while also being
the \emph{fastest} to train.

% =====================================================================
\section{Cross-task Discussion}

Across the three tasks a consistent theme emerges: \emph{inductive
bias substitutes for data, and pretraining substitutes for both}.
CycleGAN's cycle-consistency constraint (Task~1) is a strong inductive
bias that lets it learn unpaired translations on only a few thousand
images per domain. The vanilla Transformer MT model (Task~2) has weak
inductive bias and fails at 23K pairs; mBART's multilingual pretraining
rescues performance without any in-domain fine-tuning.
CNNs' locality/equivariance bias (Task~3) substitutes for the data
scale that a from-scratch ViT would require, while pretrained ViTs
recover this gap through ImageNet-21k exposure.

% =====================================================================
\section{Conclusion}

We have implemented and evaluated three foundational deep-learning
systems covering unpaired image translation, low-resource machine
translation and image classification. Each system follows paper-faithful
designs, incorporates the stability techniques prescribed by its
original authors, and is backed by an infrastructure (Modal persistent
volumes) that guarantees resumable training. Quantitative results and
qualitative analyses are consistent with the published literature: cycle
consistency enables reliable unpaired translation with modest data,
transformer MT requires either large parallel corpora or multilingual
pretraining to avoid mode collapse, and ViTs outperform CNNs on small
image datasets only when supervised pretraining is applied.

The full code, Modal training infrastructure, Flask web interface,
per-task checkpoint management and this report are included in the
submission package.

% =====================================================================
\begin{thebibliography}{99}

\bibitem{cyclegan}
Zhu, J.-Y., Park, T., Isola, P., Efros, A.A.: Unpaired image-to-image
translation using cycle-consistent adversarial networks. In:
Proceedings of the IEEE International Conference on Computer Vision
(ICCV) (2017)

\bibitem{lsgan}
Mao, X., Li, Q., Xie, H., Lau, R.Y.K., Wang, Z., Smolley, S.P.: Least
squares generative adversarial networks. In: Proceedings of ICCV (2017)

\bibitem{vaswani}
Vaswani, A., Shazeer, N., Parmar, N., Uszkoreit, J., Jones, L., Gomez,
A.N., Kaiser, L., Polosukhin, I.: Attention is all you need. In:
Advances in Neural Information Processing Systems (NeurIPS) (2017)

\bibitem{vit}
Dosovitskiy, A., Beyer, L., Kolesnikov, A., Weissenborn, D., Zhai, X.,
Unterthiner, T., Dehghani, M., Minderer, M., Heigold, G., Gelly, S.,
Uszkoreit, J., Houlsby, N.: An image is worth $16\!\times\!16$ words:
Transformers for image recognition at scale. In: ICLR (2021)

\bibitem{mbart}
Tang, Y., Tran, C., Li, X., Chen, P.-J., Goyal, N., Chaudhary, V., Gu,
J., Fan, A.: Multilingual translation with extensible multilingual
pretraining and finetuning. arXiv:2008.00401 (2020)

\bibitem{bleu}
Papineni, K., Roukos, S., Ward, T., Zhu, W.-J.: BLEU: a method for
automatic evaluation of machine translation. In: ACL (2002)

\bibitem{bpe}
Sennrich, R., Haddow, B., Birch, A.: Neural machine translation of
rare words with subword units. In: ACL (2016)

\bibitem{eda}
Wei, J., Zou, K.: EDA: Easy data augmentation techniques for boosting
performance on text classification tasks. In: EMNLP-IJCNLP (2019)

\bibitem{cifar10}
Krizhevsky, A.: Learning multiple layers of features from tiny images.
Technical Report, University of Toronto (2009)

\bibitem{timm}
Wightman, R.: PyTorch Image Models.
\url{https://github.com/rwightman/pytorch-image-models} (2019)

\bibitem{personfacesketches}
Almightyj: Person Face Sketches dataset.
\url{https://www.kaggle.com/datasets/almightyj/person-face-sketches}.
Accessed April 2026.

\bibitem{zainparallel}
Zainuddin: Parallel corpus for English--Urdu language.
\url{https://www.kaggle.com/datasets/zainuddin123/parallel-corpus-for-english-urdu-language}.
Accessed April 2026.

\end{thebibliography}

\end{document}
